\chapter{Literature Review}

\section{Introduction}
Real estate valuation and market analysis have long been subjects of academic and practical interest. This chapter surveys the evolution of methodologies in this domain, ranging from traditional econometric approaches to modern artificial intelligence applications. It also examines the specific literature regarding the Singapore housing market to contextualize the unique challenges addressed by this project.

\section{Traditional Real Estate Valuation Models}
The foundational framework for property valuation is the \textbf{Hedonic Pricing Model}, popularized by Rosen \cite{rosen1974hedonic}. This theory posits that a heterogeneous good, such as a house, can be valued as the sum of its constituent attributes (e.g., location, size, number of bedrooms).
\begin{itemize}
    \item \textbf{Strengths:} These models provide interpretability, allowing analysts to quantify the marginal value of specific features (e.g., the premium for an extra bathroom).
    \item \textbf{Limitations:} They often rely on strict assumptions of linearity and independence among variables. Furthermore, they struggle to incorporate high-dimensional or unstructured data, such as textual property descriptions or complex spatial relationships.
\end{itemize}

\section{Machine Learning in Real Estate}
In recent years, the limitations of traditional models have led researchers to explore Machine Learning (ML) techniques.
\begin{itemize}
    \item \textbf{Automated Valuation Models (AVMs):} Advanced regression techniques (e.g., Random Forests, Gradient Boosting) have been shown to outperform linear regression in price prediction accuracy by capturing non-linear relationships \cite{park2015using}. Pow et al.\ \cite{pow2019machine} specifically demonstrated that gradient-boosted tree ensembles achieve lower MAPE than hedonic regression across multiple property segments in a UK context, a finding this project corroborates for Singapore.
    \item \textbf{Gradient Boosting (XGBoost):} Chen and Guestrin \cite{chen2016xgboost} introduced XGBoost as a scalable, regularised gradient boosting framework that has since become a standard benchmark in tabular data competitions and applied ML. Its use of second-order gradient statistics, column subsampling, and L1/L2 regularisation makes it particularly robust to overfitting on moderate-sized datasets such as the per-segment training sets (n\,=\,277--16,728) used in this project.
    \item \textbf{Model Interpretability (SHAP):} Lundberg and Lee \cite{lundberg2017unified} introduced SHAP (SHapley Additive exPlanations), which provides theoretically grounded, additive feature attributions for any model. Unlike earlier importance measures (e.g., Gini impurity), SHAP values satisfy consistency and local accuracy axioms, making them suitable for communicating property-specific valuation drivers to non-technical end users.
    \item \textbf{Natural Language Processing (NLP):} Recent studies, such as Baur et al.\ \cite{baur2022automated}, utilise NLP to extract sentiment and semantic features from listing descriptions, demonstrating that ``soft'' factors significantly influence pricing. The present project extends this direction by using an LLM (Claude) not for feature extraction but for \textit{query understanding}: translating free-text search intent into structured database filters.
    \item \textbf{Neural Networks:} Deep learning architectures allow for the integration of multi-modal data (text and images), providing a more holistic view of property value \cite{bian2020comparing}. This project prioritises interpretability and data efficiency over raw predictive power, hence gradient-boosted models are preferred over neural networks for the valuation task.
\end{itemize}

\section{The Singapore Context: A Dual-Tier Market}
Research on the Singapore market highlights its unique dual structure. Tu \cite{tu2004dynamics} analyzed the dynamics of the Singapore private housing market, finding significant relationships between private housing prices and the resale public housing sector. Sing et al. \cite{sing2006price} further demonstrated price dynamics between public and private housing, where rising private prices can pull up HDB resale prices due to the substitution effect and inter-market mobility.

\section{Research Gap}
While existing literature covers individual components—valuation algorithms, market segmentation, or NLP feature extraction—three gaps remain unaddressed in the Singapore context.

First, no prior published system aggregates and deduplicates property listings \textit{across all four major Singapore portals} (PropertyGuru, 99.co, EdgeProp, SRX) into a single unified dataset. Most local studies, including Bian et al.\ \cite{bian2020comparing}, rely on a single structured transaction dataset (e.g., URA REALIS), which lacks the breadth and recency of live listing data.

Second, existing ML valuation work on Singapore property \cite{bian2020comparing} does not provide per-prediction SHAP explanations accessible to end users. The opacity of black-box models limits their practical utility for buyers and agents who need to understand \textit{why} a valuation estimate is high or low.

Third, no prior work applies LLM-based natural language interfaces to Singapore property search. Users of existing portals are constrained to rigid filter menus that cannot express nuanced requirements such as ``expat family, near international school, under S\$8,000/month.'' This project bridges all three gaps by building an integrated platform that unifies multi-portal data, surfaces SHAP-grounded valuations, and enables conversational search with agentic fallback behaviours.

To the best of the author's knowledge, no prior published system simultaneously addresses all three gaps: multi-portal Singapore property aggregation with composite-key deduplication, per-prediction SHAP-grounded valuation explanations, and an LLM-backed agentic search pipeline with progressive filter relaxation and a grounded valuation chat assistant---within a single end-to-end deployable platform.
